\section{Direct comparisons at simultaneous births}\label{sec:direct}

We prove the part of Proposition~\ref{tab:first-join-input} concerning
adjacent simultaneous births. The main distinction is between an exact coincidence
after replacing phases by recognized rationals, which is charged to first
mass, and a nonexact coincidence, for which the arithmetic extraction
estimate supplies a summable saving. An additional birth discard will
exclude the rational concentrations arising in one of the extracted cases.
It depends only on the original lists of good assigned rows.

\subsection{The comparison statement and determinant coordinates}

A \emph{direct pair} is an ordered pair of centres inserted at the same
height $Y=Y_0$ and satisfying $\log UT\le\beta Y$, with self-pairs allowed.
Simultaneous newborn pairs that first share a component but satisfy
$\log UT>\beta Y$ are treated in Section~\ref{sec:joins}.
On medium type, the band separation implies $B_i=B_j$; moreover, the two
centres have the same product $P$ and the same tagged/tail division, since
$J\gg\beta D_0$. On high type they have the same last dense height $Y_d$
and the same $P$. Indeed, any graph neighbour of a high centre at its
insertion has the same last dense bin, so its entire component is newborn.
Group centres by this common $P$, start, and tagged/tail division, including
the band on medium type but not splitting bands on high type. Each row
belongs to one group, and direct pairs do not cross groups. Write
\[
 X_0=B_i=B_j\quad\hbox{on medium type},
 \qquad X_0=Y_d\quad\hbox{on high type}.
\]
The definitions of the starts and the centre-scale bounds give
\begin{equation}\label{eq:direct-size}
 \log\bigl(2+UT+K+r_{\max}/r_{\min}\bigr)\ll X_0.
\end{equation}

For a centre $i$ and an outer mask $a\mid P_I$, let $\nu_{i,a}$ be its
labelled birth point measure with the raw static weights, including the
soft tests and the additional discard constructed below, but without the
coefficient $\pi_{P_I}(a)$ and before any adaptive update or cap clipping.
Thus
\[
 \lambda_i=\sum_{a\mid P_I}\pi_{P_I}(a)\nu_{i,a}.
\]
Point-pair sums below are per unit length: the first location is counted in
one unit interval and the second ranges over the periodic lift. Products
of anchor-averaged weights are evaluated with independent anchor choices,
also for self-pairs and coinciding labels.

\begin{proposition}[Direct birth bounds]\label{dir:birth-comparisons}
Use the retained centres of one type and colour, the good masks, and the
birth prescriptions of Sections~\ref{sec:centres}
and~\ref{sec:prescriptions}, with the hierarchy fixed in
Section~\ref{cen:hierarchy}. Let the finite block have mass $W\le1$.
For every fixed $C>0$, the sum over all
direct pairs of
\[
 r_{\min}\iint_{|x-y|\le C r_{\max}}
       d\lambda_i(x)\,d\lambda_j(y)
\]
is $O(1)$. The same-width, one-law sum
\[
 \sum_{\substack{(i,j)\text{ direct}\\r_i=r_j=r}}
 r\sum_{a\mid P_I}\pi_{P_I}(a)
 \iint_{|x-y|\le C r e^{-E_gX_0}}
       d\nu_{i,a}(x)\,d\nu_{j,a}(y)
\]
is $O(W)+O_{k_0}(W^{1+c})$ for some fixed $c>0$. Each ordered centre pair
is counted once, at its common birth. Constants may depend on the fixed
hierarchy, cutoffs, and distance constant, but not on the finite block or
the slot cap $N_0$.

The additional discard used here is a static rule on medium groups,
determined by their good assigned-row lists before any running-table
deletions. Given any fixed $\delta>0$, its total first-mass loss is at most
$\delta W$ after choosing the fixed threshold $\lambda_0$ sufficiently
large. It uses only squarefree low numerator moduli. The comparison bounds
remain valid after any further nonincreasing thinning of birth weights.
\end{proposition}

We first establish the estimates common to both sums. Write
$L_i=S_iP$, $L_j=S_jP$, and denote the low row denominators by $V,W'$.
Set
\[
 s=\frac{W'}{(V,W')},\qquad t=\frac{V}{(V,W')}.
\]
Let $m_i,m_j$ be the tagged masks and let
$\alpha_i=a_i\gamma$, $\alpha_j=a_j\gamma$ be the outer phases. In the
narrow sum $a_i=a_j$ with one law; in the broad sum they have independent
laws. The processed masks and the integers used by the tag prescriptions
are
\[
 d_i=(S_i/V)m_i,\qquad d_j=(S_j/W')m_j,
 \qquad \mathcal L_i=\lfloor d_i\alpha_i\rfloor,
 \quad \mathcal L_j=\lfloor d_j\alpha_j\rfloor.
\]
Define
\[
 f=(m_i,m_j),\qquad f_0=(Td_i,Ud_j),
 \qquad \sigma_0=\frac{Td_i}{f_0}=\frac{s m_i}{f},
 \quad \tau_0=\frac{Ud_j}{f_0}=\frac{t m_j}{f}.
\]
The identities follow from the disjoint prime supports of low and tagged
parts; in particular, $\sigma_0$ and $\tau_0$ are coprime. A pair within
the tested distance, henceforth a \emph{cohit}, satisfies
\begin{equation}\label{eq:direct-determinant}
 \begin{split}
 |\sigma_0(z_i+\alpha_i)-\tau_0(z_j+\alpha_j)|&\le\Delta,\\
 \Delta&=\frac{CK\varepsilon}{f_0}
        =\frac{C[V,W']P r_{\max}\varepsilon}{f},\\
 \frac{\Delta}{\sigma_0}
        &=C\psi_i\varepsilon\frac{r_{\max}}{r_i},
 \qquad \psi_i=\frac{VP r_i}{m_i},
 \end{split}
\end{equation}
where $\varepsilon=1$ for broad comparisons and
$\varepsilon=e^{-E_gX_0}$ for narrow ones. Constants in distance tests are
absorbed into $C$. For each determinant integer the bare count, multiplied
by $r_{\min}$, is at most
\[
 \Theta_{ij}\frac{f_0}{d_i d_j}.
\]
Indeed, writing $G=(L_i,L_j)$, the projected denominators
$L_i/d_i,L_j/d_j$ have gcd
$Gf_0/(d_i d_j)$, and
\[
 Td_i=[S_i,S_j]m_i/V,\qquad
 Ud_j=[S_i,S_j]m_j/W',\qquad
 f_0=[S_i,S_j]f/[V,W'].
\]
Put $\eta_1=[V,W']P r_{\max}$. Call a comparison \emph{ultra} if it is
narrow, or if it is broad and either $f>1$ or
$-\log\eta_1\ge Q_0X_0$. The tag-prime lower bounds, the choices of
$J,E_g$, and \eqref{eq:direct-size} imply
\[
 \Delta\ll e^{-.9Q_0X_0}
 \qquad\text{in every ultra comparison}.
\]

Every narrow comparison is therefore ultra. Lemmas~\ref{dir:ultra-switch}
and~\ref{dir:tag-fold} reduce this regime to rational models and low
coordinates; Lemmas~\ref{dir:exact-equalities} and~\ref{dir:ultra-count}
then complete its count. The remaining broad comparisons have $f=1$
and are controlled by the row kernel of Lemma~\ref{dir:row-kernel}.
Before extracting a small-scale box, we fix the static holes in
Construction~\ref{dir:hole-rule}. The final three lemmas exclude high,
shallow medium, and deeper medium boxes, respectively.

\subsection{Ultra switches and tag folding}

The first lemma extracts rational phases at ultra tolerance. A switch is
only a change of summation parameters for an upper bound; the changed
factor is not required to produce another assigned row or another cohit.

\begin{lemma}[Ultra switching]\label{dir:ultra-switch}
For ultra comparisons, discard an exceptional count of total size
$O_{k_0}(W^{1+c})$, for some $c>0$. On the remaining pairs, apply all
divisor switches
\[
 h\mid(d_i,\sigma_0),\qquad
 e^{C_1D'X_0}\lesssim h\lesssim e^{3C_1D'X_0},
\]
all prime switches above the lower cutoff, and the corresponding switches
on side $j$. If any switch is available, there are coupled rationals
$A_i,A_j$ of denominators at most $e^{O(D')X_0}$ such that
\[
 |A_i-\sigma_0\alpha_i|+|A_j-\tau_0\alpha_j|
       \ll\Delta e^{O(D')X_0},
 \qquad A_i-A_j=-(\sigma_0z_i-\tau_0z_j).
\]
They are unique within these bounds. For every switched divisor $h$ on
side $i$, the reduced denominator of $A_i/h$ is at most
$e^{C_3D'X_0}$, where $C_3\ll C_1$ is absolute; the analogous assertion
holds on side $j$. In particular, switched large primes divide the
appropriate reduced numerator. There is a fixed $E'\ll E$ such that
$A_i/\sigma_0$ and $A_j/\tau_0$ have denominators at most $e^{E'X_0}$.
If no switch is available, both $\sigma_0,\tau_0\le e^{E'X_0}$.
\end{lemma}

\begin{proof}
Every $h\mid(d_i,\sigma_0)$ can be removed from $d_i$ without changing
$f_0$: at each affected prime the first argument of the gcd has an exponent
surplus at least the removed exponent. Fix the centres, widths, outer
masks, $d_i/h$, and $d_j$. On a dyad for $h$, equation
\eqref{eq:direct-determinant} is a rotation condition with $h$ occurring
linearly, fixed slope $\sigma_0\alpha_i/h$, and unchanged tolerance.
Lemma~\ref{ar:rotation}, with budget $D'X_0$, therefore gives the
rational alternative except on $O(e^{-D'X_0})$ of that dyad's normalized
count. Its approximation error for the slope is at most
$\Delta e^{C_3D'X_0}/h$.

Stripping $h$ contributes exactly $1/h$ to the bare ratio
$f_0/(d_i d_j)$. We may drop all other assignment and eligibility
restrictions when summing the exception. At equal centre exponents the
local unrestricted sum is at most
\[
 \sum_{a,b\ge0}p^{-\max(a,b)}=1+O(1/p).
\]
At unequal exponents, say $\nu_p(U)=g>0$, it is at most
\[
 \sum_{a,b\ge0}p^{\min(a,g+b)-a-b}
       \le (1+g)^C(1+C/p)
\]
with absolute $C$. The processed-prime harmonic sum is $O(Y_0)$, so the
stripped-mask sum costs at most
$\exp(CY_0)\tau(UT)^C$. Outer masks still integrate as probabilities,
including the one-law coupling in the narrow case. There are at most
$\exp(O(Y_0))$ dyads in all the indicated tests. Choosing $D'$ large
against these costs bounds the exceptional count for a centre pair by
\[
 \Theta_{ij}e^{-D'X_0/2}.
\]
The centre kernel \eqref{eq:centre-kernel}, lower bounds
\eqref{eq:centre-lower}, and \eqref{eq:direct-size} sum these errors to
$O_{k_0}(W^{1+c})$, leaving exponential decay for shell multiplicities.

For precision at the switch endpoints, round the lower cutoff up to a
dyad start, use divisor dyads from there through its square, and use prime
dyads from there onwards. On every nonexceptional pair every available
test gives its rational alternative. Multiplication by $h$ gives $A_i$;
the determinant relation defines
$A_j=A_i+\sigma_0z_i-\tau_0z_j$. Candidates obtained from any two
switches have denominators at most $e^{C_3D'X_0}$ and differ by at most
$O(\Delta e^{C_3D'X_0})$. The latter is smaller than the reciprocal
product of their denominators, by $Q_0\gg D'$. Thus the candidates
coincide. This proves the exact displayed relation and simultaneous
uniqueness, and $A_i/h$ has the stated small denominator for every
switched $h$.

It remains to divide by the whole signature. We have
$\sigma_0/(d_i,\sigma_0)\mid T$: at a prime, writing
$a=\nu_p(d_i)$, $b=\nu_p(T)$, $c=\nu_p(f_0)$, the exponent in this
quotient is $\max(b-c,0)\le b$. The good-mask condition bounds the
prime-power excess on $(d_i,\sigma_0)$ by $e^{C_{\rm mask}B_i}$, and every
sufficiently large prime there divides the reduced numerator of $A_i$.
Among the smaller distinct primes not dividing that numerator, the product
is below the lower switch cutoff, up to a fixed factor. Otherwise, a
first-crossing subproduct would lie between the cutoff and its square and
would be a tested divisor. Dividing by this product, which is coprime to
the numerator, would give a denominator larger than the allowed
$e^{C_3D'X_0}$. Explicitly, if $A_i=n/q$ is reduced and $(n,h)=1$,
then the reduced denominator of $A_i/h$ is $qh$, even if $h$ and $q$
share primes. The remaining factor $T$, excess powers, and subcutoff
product are all singly exponential in $X_0$ by
\eqref{eq:direct-size}; this proves the denominator bound with a fixed
$E'\ll E$. If no tested divisor occurs, the same factor decomposition
bounds the signatures themselves. The argument on side $j$ is identical.
\end{proof}

We next average the high numerator coordinates. The exact lift used in
this averaging is
\[
 D=T(d_i z_i+\mathcal L_i)-U(d_j z_j+\mathcal L_j).
\]
It is invariant under joint unit translation, and every cohit satisfies
\[
 |D|\ll T+U+K.
\]
The bound follows by replacing the two floors with their arguments and
using \eqref{eq:direct-determinant}.

\begin{lemma}[Tag arrangements and folding]\label{dir:tag-fold}
Apart from the exceptions of Lemma~\ref{dir:ultra-switch}, tag masks in
each bin may be restricted to the following arrangements: they are
identical, or they are disjoint. There are only boundedly many bin
arrangements. If $f=1$, all bins are disjoint.

For a fixed arrangement, let $P_{\rm eq}$ and $P_{\rm dis}$ be the products
of full prime powers in its identical and disjoint bins. Under
\[
 y=(P/f)x,\qquad \rho_i=(P/f)r_i,\qquad m_i'=m_i/f,
\]
the folded low arrays have the form
\[
 y_i=(Z_i+m_i'\alpha_i)/V,
 \qquad Z_i\in\mathbb Z,
\]
with transformed low tests. An upper bound for their tag-averaged
coefficient is, up to a fixed constant, the prescribed outer-law
coefficient times
\[
 f\pi_{P_{\rm eq}}(f)
       \pi_{P_{\rm dis}}(m_i')\pi_{P_{\rm dis}}(m_j').
\]
The low arrays may be counted without the original integrality
$m_i'\mid Z_i$, which has already been included in this coefficient.
One may still retain the predicate that a surviving actual cohit exists
at the original masks and anchor choices.
\end{lemma}

\begin{proof}
First suppose a bin's masks are disjoint, and fix both anchor choices.
At a prime full on both sides, two active forbidden residues in their
common high-coordinate translation can coincide only if
$D\equiv TF_i-UF_j\pmod p$. The size bound for $D$, the reconstruction
height and denominator bounds, and the bin's prime size turn this
congruence into equality over $\mathbb Q$. At an anchor of side $i$,
absent from side $j$'s mask, reduction then gives
$d_j z_j+\mathcal L_j\equiv F_j$, contradicting the hard test. When one
side has no key, its scalar $R_p$ already supplies the required product
factor and no coincidence argument is needed.

If $f>1$, the pair is ultra and every mismatched tag prime is among the
switches of Lemma~\ref{dir:ultra-switch}. If a mismatch exists, define
\[
 F_i^{\rm pred}=\mathcal L_i-\frac{S_i}{[V,W']}f A_i,
 \qquad
 F_j^{\rm pred}=\mathcal L_j-\frac{S_j}{[V,W']}f A_j.
\]
The denominator of the low-centre ratio on side $i$ divides $T$, and the
one on side $j$ divides $U$. More importantly,
\[
 \frac{S_i}{[V,W']}f\sigma_0\alpha_i=d_i\alpha_i.
\]
Thus the potentially large phase cancels against the floor, giving size
$O(1+e^{O(D')X_0}L_i r_{\max})$ and the corresponding small denominator.
These values fit the tag reconstruction cutoffs. At a shared omission the
multiplier contains the anchor prime; at a mismatched omission numerator
divisibility comes from its switch. Every anchor in a nonempty bin
therefore reconstructs the predicted tag. The exact determinant relation
gives
\[
 T F_i^{\rm pred}-U F_j^{\rm pred}=D.
\]
In a bin where both masks are nonempty, reduction at a mismatch prime
violates the full side's forbidden-residue test. Such bins consequently
have identical masks; bins with an empty mask are disjoint. If there was
no mismatch, all masks were already identical. The number of tag bins is
bounded by the fixed hierarchy, so the number of arrangements is bounded.

For the fold, the original integrality is $Z_i=m_i'z_i$. Count first over
a $y$-interval of integer length $P/f$, using the generic arrays and
averaging joint shifts
\[
 Z_i\longmapsto Z_i+Vk,\qquad
 Z_j\longmapsto Z_j+W'k,\qquad k\bmod P/f.
\]
Replacing $r$ by $\rho$ is exactly the change from the original unit
length to the folded unit length. In a disjoint bin, a one-sided omission
has integrality density $1/p$. Relative to the two deficit-law factors
this costs at most an extra fixed power of $R_p^{-1}$; the product of
these losses is bounded by the good-mask harmonic condition. At a prime
omitted by neither side, the preceding noncoincidence argument gives
upper mean $R_p^2$. The slopes are units: in these coordinates
$d_i z_i=(S_i/V)fZ_i$, and its translation slope is $S_i f$. If a
coincidence occurs, the lifted identity forbids survival on every
translation satisfying integrality.

In an identical bin, keep just one factor $R_p$ at a full prime. A shared
omission has one residual scalar $R_p$ when the full centre exponent is at
least two, and no factor when it is one. These are precisely the local
factors of $f\pi_{P_{\rm eq}}(f)$. More explicitly, for full centre
exponent $\nu$, the local upper bounds are
\[
\begin{array}{c|c|c}
 \text{arrangement and deficits}&\nu=1&\nu\ge2\\ \hline
 \text{disjoint},\ (0,0)&R_p^2&R_p^2\\
 \text{disjoint},\ (1,0)&1/p&R_p/p\\
 \text{identical},\ (0,0)&R_p&R_p\\
 \text{identical},\ (1,1)&1&R_p
\end{array}
\]
with the symmetric one-sided case understood. The one-sided row differs
from $\pi_{p^\nu}(p)\pi_{p^\nu}(1)$ by at most $R_p^{-1}$; the last
row is $p\pi_{p^\nu}(p)$. All required digits occur in
$k\bmod P/f$; in particular a shared omitted prime of exponent one needs
no further digit. The digit CRT multiplies these upper means. Low tests
use $z_i=Z_i/m_i'$ modulo low powers, with division by a unit, and are
unchanged by the high translations. This proves the folded coefficient.
Since all majorizations are nonnegative, an existence predicate for an
actual cohit may be retained outside the count. This predicate is only
the assertion that some actual integral cohit survives at the fixed
row, mask, and anchor data: if none does, the original count is zero;
otherwise the unconditional digit upper bound applies. No pointwise
survival condition is factored through the digit average.
\end{proof}

\subsection{Billing exact modeled equalities}

A cohit has an \emph{exact modeled equality} if replacing its two phases
by recognized rationals from the low birth tiers makes its two lifted
locations exactly equal. Equality here is in the line, not merely modulo
one. This case requires a first-mass estimate, since it cannot yield a
determinant saving.

\begin{lemma}[Exact equality bound]\label{dir:exact-equalities}
The total folded comparison count with exact modeled equality, using all
recognized tiers and either the broad or narrow outer-mask law, is
$O(W)$.
\end{lemma}

\begin{proof}
\emph{Independent mask averaging.}
Fix a group and one of the arrangements of Lemma~\ref{dir:tag-fold}, then
condition on $f$ and, for narrow comparisons, on the common outer mask.
Drop pairwise restrictions on centres, masks, and widths. In the
resulting upper bound, use the unrestricted independent deficit laws
for the remaining tagged masks, and also for the outer masks in the
broad case.

For each qualifying folded row and chosen model
$\zeta_i=n_i/q_i$, form a vector of functions indexed by rational
$y_*\in[0,1)$. Its coordinate at $y_*$ is zero unless
\[
 y_*=(Z_i+m_i'\zeta_i)/V,\qquad Z_i\in\mathbb Z.
\]
Otherwise that coordinate is the interval indicator supported on
$|x-y_i|\le C'\rho_i$, multiplied by the low-test weight at $Z_i$.
Choose $C'$ sufficiently large compared to the distance constant.
The inner product of two such vectors majorizes, up to a fixed factor,
the low weighted cohit count times $\min(\rho_i,\rho_j)$ when there is
exact modeled equality. Choosing the first lift by $y_*\in[0,1)$ is
equivalent to counting per unit length: joint integer translation
preserves every transformed low test.

Let $\mathbf b$ denote this vector, and include all assigned good rows
and all recognized tiers at each sampled mask. Since the remaining laws
are the same on both sides, the widened comparison is bounded by
\[
 \left\|\mathbb E\sum_{\rm rows,tiers}\mathbf b\right\|^2
 \le
 \mathbb E\left\|\sum_{\rm rows,tiers}\mathbf b\right\|^2.
\]
The fixed outside coefficient is the one in
Lemma~\ref{dir:tag-fold}. This is an ordinary Hilbert-space Jensen
inequality; the two masks are coupled only in the expression on the
right. The common unrestricted law is
$\pi_{P_{\rm dis}}(m')\pi_{P_I}(a)$ in the broad case and
$\pi_{P_{\rm dis}}(m')$ after fixing the common outer mask in the narrow
case. Assignment and good-mask indicators remain inside the vector,
as zero contributions from unavailable rows; they do not condition
this law. Split the bounded number of tiers by Cauchy--Schwarz,
distinguishing the large broad, small broad, and narrow parameter
sets on medium type.

\emph{Common-model partner counts.}
We may therefore fix one common mask $m$, one outer phase $\alpha$, and
one tier. Within this tier, the same reduced rational $\zeta=n_0/q$
models $\alpha$ wherever a model is recognized, even across widths, by
Lemma~\ref{pre:model-uniqueness}. Moreover $(q,m')=1$, where
$m'=m/f$, since the model denominator is below the tag-prime cutoff.
At these masks a specified $V$ determines the actual denominator
$VP/(ma)$, and hence determines at most one assigned row, including its
width and centre.

For an exact equality between $V$ and $W'$, write
\[
 H_V=2+q+b_V,\qquad
 b_V=\frac{|\alpha-\zeta|m}{VP r_i},
\]
and define $H_{W'},b_{W'}$ similarly. The soft tests and the identity
\[
 s(qZ_i+m'n_0)=t(qZ_j+m'n_0)
\]
imply
\begin{equation}\label{dir:soft-ratios}
 q\mid s-t,\qquad
 s\le H_{W'}^{h_s},\qquad t\le H_V^{h_s}.
\end{equation}
Indeed, $(s,t)=1$ gives $t\mid qZ_i+m'n_0$ and
$s\mid qZ_j+m'n_0$. Since $t\mid V$, division by the unit $m'$ in the
low test bounds $t$ by its soft gcd; the other side is the same.
Reduction modulo $q$, using $(q,m'n_0)=1$, gives the first assertion.
If $\alpha\ne\zeta$ and the two interval bumps intersect, their centres
also give
\begin{equation}\label{dir:bias-overlap}
 m'|\alpha-\zeta|\,|1/V-1/W'|\ll\rho_i+\rho_j,
 \qquad
 |s-t|\ll s/b_V+t/b_{W'}.
\end{equation}

Orient so that $H_V\le H_{W'}$. First suppose
$H_{W'}\le H_V^{1/(3h_s)}$. Then
$s\le H_V^{1/3}$ and $t\le H_V^{h_s}$. For large $H_V$, either
$q>s+t$, in which case \eqref{dir:soft-ratios} forces $s=t$, or
$q\le s+t\le2H_V^{1/3}=o(H_V)$. In the latter case
$b_V\ge H_V/2$ and $b_{W'}\ge H_{W'}/2$ for large $H_V$, so the
right side of \eqref{dir:bias-overlap} is
$O(H_V^{-2/3}+H_V^{h_s-1})<1$ and again $s=t$.
The zero-bias case belongs to the first alternative at large height.
Thus large comparable heights give the same row, while bounded heights
permit only boundedly many pairs of ratios $(s,t)$.

Now suppose $H_{W'}>H_V^{1/(3h_s)}$. Then $b_{W'}\asymp H_{W'}$,
and
\begin{equation}\label{dir:disparate-widths}
 \frac{\rho_j}{\rho_i}
   =\frac{V}{W'}\frac{b_V}{b_{W'}}
   \ll\frac{H_V^{1+h_s}}{H_{W'}}.
\end{equation}
The identity is needed only when the bias is nonzero; zero bias makes
$H_V=H_{W'}$ and cannot occur in this case. Per point on $V$, the number
of possible partners in a dyad $H_{W'}\asymp D$ is
$O(H_V^{h_s}D^{h_s})$, by \eqref{dir:soft-ratios}. Each ratio determines
$W'$, which determines its row, and exact equality determines its point.
Consequently the sum of partner lengths relative to $\rho_i$ is at most
\[
 O\left(
 H_V^{1+2h_s}
 \sum_{D\gtrsim H_V^{1/(3h_s)}}D^{h_s-1}
 \right)=O(1).
\]
Here the sum is dyadic, and
$1+2h_s+(h_s-1)/(3h_s)<0$ for $h_s<.01$. Since $H_V\ge3$, this is
uniform also for bounded $H_V$ and infinitely many possible partner
height dyads. In the comparable case,
charge the first bump length and use the bounded number of partners.
Products of low weights are at most the charged row's low weight.
Its average is at most $R(V)$, so the same-mask square is bounded by
constant times low first mass.

\emph{Return to assigned first mass.}
We have
\[
 f\pi_{P_{\rm eq}}(f)\rho_i
       =P\pi_{P_{\rm eq}}(f)r_i.
\]
Averaging over $f$, the remaining tagged masks, and the outer masks
returns the original row currency, because
\[
 \begin{aligned}
 \pi_{P_{\rm eq}}(f)\pi_{P_{\rm dis}}(m')&=\pi_{P_M}(m),\\
 \phi(v)&=\phi(V)\phi(P_M/m)\phi(P_I/a)\\
        &=P\phi(V)\pi_{P_M}(m)\pi_{P_I}(a),
 \end{aligned}
 \qquad m=fm',
 \quad v=VP/(ma).
\]
In particular,
\[
 r_iP\phi(V)\pi_{P_M}(m)\pi_{P_I}(a)=r_i\phi(v).
\]
The bounded number of arrangements and tiers only changes the constant.
\end{proof}

\begin{lemma}[Completion of ultra comparisons]\label{dir:ultra-count}
All ultra comparisons in Proposition~\ref{dir:birth-comparisons} have
total cost $O(W)+O_{k_0}(W^{1+c})$.
\end{lemma}

\begin{proof}
When a switch is available, division of the approximants of
Lemma~\ref{dir:ultra-switch} by $\sigma_0,\tau_0$ gives phase errors,
in units of $\psi_i,\psi_j$, bounded by
$\varepsilon e^{O(D')X_0}$; here
\eqref{eq:direct-size} absorbs the width ratio. These rationals qualify
in the birth tiers: the large broad tier in a medium broad comparison
is enabled by $\psi_i,\psi_j\ll\Delta$; a medium narrow comparison uses
the narrow tier; a high comparison uses its unique broad tier.
The choices $E\gg E'$, $Q_0\gg E$, and $E_g\gg E$ ensure the required
denominator, height, and error bounds. Their coupled determinant
identity gives exact modeled equality, covered by
Lemma~\ref{dir:exact-equalities}.

It remains to count nonexact pairs with no switch. Fix the group and the
two outer and tagged masks, treating broad and narrow comparisons
separately. There can be at most one coprime signature vector
$(\sigma_0,\tau_0)$ among these pairs. Indeed, two distinct coprime
vectors give an invertible pair of equations
\eqref{eq:direct-determinant}. Solving them with the integral
determinants as exact right sides gives rational replacements of the
two fixed phases, of denominators at most $2e^{2E'X_0}$, and errors
at most $e^{O(E')X_0}$ times the larger of the two tolerances.
At the pair with that larger tolerance, each error is at most the
corresponding numerator width, $\psi_i$ or $\psi_j$, times
$\varepsilon e^{O(E')X_0}$. Explicitly, put $M=e^{E'X_0}$.
The coefficient determinant is a nonzero integer of absolute value
at most $2M^2$, so the phase errors are at most $2M\Delta_{\max}$.
On the pair attaining $\Delta_{\max}$, both sides satisfy
\[
 \frac{\Delta_{\max}}{\psi_i}
      =C\varepsilon\sigma_0\frac{r_{\max}}{r_i}
      \le C\varepsilon M e^{CX_0},
\]
and the analogous inequality with $\tau_0,\psi_j$.
The same tier verification as above would make
that pair exact modeled, a contradiction.

For the one remaining signature, a fixed $V$ has at most one partner
$W'$. The folded determinant integer is unique, since the tolerance is
ultra. Retain only one low factor at each prime of $V$ not dividing $t$,
and average over its full common digit in the determinant cycle. The
resulting count is at most
\[
 \frac{V}{t}R(V)\prod_{p\mid t}R_p^{-1}
       =\frac{\phi(V)}{\phi(t)}
       \le\phi(V).
\]
Charging $\rho_i$ times this bound to the row on $V$ and restoring the
fold coefficients as in Lemma~\ref{dir:exact-equalities} gives $O(W)$.
This one-neighbour argument includes all widths at the fixed masks.
The ultra-switch exceptions supply the remaining
$O_{k_0}(W^{1+c})$ term.
\end{proof}

\subsection{The non-ultra row kernel}

All narrow comparisons have now been counted. For the remaining broad
comparisons, $f=1$. Use folded widths
$r_i'=Pr_i$, $r_j'=Pr_j$, width ratio
$\mathcal R=r_{\max}/r_{\min}$, and phases
$m_i\alpha_i,m_j\alpha_j$. The \emph{primary model} is the recognized
broad model, if present, and otherwise the recognized narrow model.
If neither is recognized, there is no primary model and all retained
low factors are ordinary scalars $R_p$. When a primary model exists,
retain its hard tests, with an ordinary scalar $R_p$ wherever it has no
active exclusion. This upper version costs at most a fixed multiplicative
constant even if two distinct medium models were recognized. In that
case the extra hard primes of the narrow model are above $e^{cB}$,
whereas low primes are at most $e^{JB}$, so their reciprocal sum is
bounded. Exact primary-modeled equalities have already been counted.

The homogeneous antecedent of this determinant-sieving strategy is the
Pollington--Vaughan overlap estimate~\cite{PV1990}, as restated in
\cite[Lemma~5.3]{KM2020}. The shifted, weighted row kernel with the
model-dependent restrictions used here is proved below; it is not
asserted to follow from that homogeneous estimate.

\begin{lemma}[Nonexact row kernel]\label{dir:row-kernel}
At fixed masks, the weighted folded count of the remaining broad
determinants is at most a constant times
\begin{equation}\label{eq:row-kernel}
 (r_i'\phi(V))(r_j'\phi(W'))
 \max(1,1/\eta_1)
 \prod_{\substack{p\mid st\\p>\eta_1^{c_1}}}R_p^{-C},
\end{equation}
where $c_1>0$ is a small absolute constant and mask probabilities are
applied outside the expression. This bound may retain the predicate
that there is an actual integral, surviving, nonexact cohit for the
specified masks and anchors, including its hole tests.
\end{lemma}

\begin{proof}
At a common prime not dividing $st$, the determinant cycle traverses a
full common digit. If two active forbidden residues coincide, the prime
divides the nonzero integer
\[
 q_iq_j\bigl\{s(Z_i+m_i\zeta_i)
                 -t(Z_j+m_j\zeta_j)\bigr\},
\]
where $q_i,q_j$ are the primary-model denominators. Each phase bias
contributes at most a constant times its relative height times
$\eta_1$, and $q_i,q_j$ are bounded by those heights. Thus this integer
has absolute value
$O(\eta_1\max(H_v,H_{\widehat v})^3)$, where $v,\widehat v$ are the
actual rows associated with the folded denominators $V,W'$.
Both hard restrictions require
$p>\max(H_v,H_{\widehat v})^{\epsilon_s}$. Above
$\eta_1^{c_1}$ only boundedly many prime divisors can satisfy both
thresholds. Indeed, writing $H=\max(H_v,H_{\widehat v})\ge3$ and
$M=\max(2,\eta_1^{c_1},H^{\epsilon_s})$,
\[
 \log(C\eta_1H^3)
 \le \log C+\log^+\eta_1+3\log H
 \ll_{C,c_1,\epsilon_s}\log M.
\]
The integer is nonzero precisely because this lifted determinant is not
an exact primary-modeled equality. If its upper size bound is below
one, no such integer occurs at all. Thus the coincidence loss is
uniform even when $\eta_1<1$ or $H$ is large. If either factor is
a scalar, there is no coincidence loss. At mismatch primes, pay the
displayed product.

When $\eta_1$ is large, average larger common digits first, keeping the
smaller ones fixed, then widen the determinant set and apply
Lemma~\ref{ar:sieve} through $\eta_1^{c_1}$. For bounded $\eta_1$ the
remaining small primes cost only a constant. Write $g=(V,W')$ and
$r_{\min}'=\min(r_i',r_j')$. The resulting width-weighted bound is
$g r_{\min}'\max(1,\eta_1)R(V)R(W')$ times the displayed mismatch
product. Since $\eta_1=[V,W']r_{\max}'$, its mass factor is
\[
 g r_{\min}'\max(1,\eta_1)R(V)R(W')
 =(r_i'\phi(V))(r_j'\phi(W'))\max(1,1/\eta_1).
\]
This proves \eqref{eq:row-kernel}. A nonnegative count vanishes when its actual-cohit
predicate fails, so the predicate may be retained throughout the
majorization.
\end{proof}

\paragraph{A preliminary linear charge.}
For small determinant scales use bins
$e^{-\lambda}\le\eta_1<e^{1-\lambda}$ with integer
$\lambda\ge\lambda_0$, and put $\eta=e^{-\lambda}$.
Before aggregating the broad comparisons in the medium range
$\ell_{\rm sh}B\le\lambda\le Q_0B+O(1)$, charge nonexact pairs for which
\begin{equation}\label{dir:small-signatures}
 s m_i,\ t m_j,\ \mathcal R\le e^{h'B},
\end{equation}
where the fixed exponent $h'$ is sufficiently small against $h_0$.
At fixed outer and tagged masks there is at most one such coprime
signature over the whole range. Otherwise the two-equation argument
of Lemma~\ref{dir:ultra-count} gives exact phase rationals of small
denominator, with relative errors at most $e^{O(h')B}$ in units of
the widths at the pair with the larger tolerance. At that pair
$\psi_i,\psi_j\ll e^{-\ell_{\rm sh}B}$, so the small broad tier, or
the larger tier, is enabled. The solved models qualify and force
exact modeled equality, a contradiction. The one-neighbour count
from Lemma~\ref{dir:ultra-count} therefore gives a linear charge
$O(W)$. We exclude this family from the remaining edge counts,
without changing the aggregate good-row lists.

For an exact width in a group, aggregate the low vertices $V$ with
normalized weights
\begin{equation}\label{dir:aggregate-vertices}
 r_i'\phi(V)\sum_{m,a}\pi_{P_M}(m)\pi_{P_I}(a)
 \mathbf1_{\{\text{the row }VP/(ma)\text{ is assigned here at width }r_i\}}.
\end{equation}
Only retained good masks enter this sum. Its total is the ordinary
block mass on the sublist. After division by $r_i'$, each weight is
at most $\phi(V)$: fixed $V,m,a$ prescribe one actual row, and the
unrestricted mask laws have total mass one. Edge weights below are
products of these normalized weights, with the mismatch factors and
necessary cohit predicates incorporated inside the mask and anchor
averages, excluding the pairs already charged in
\eqref{dir:small-signatures}. Their totals, multiplied by
$\max(1,1/\eta_1)$ on a scale,
bound Lemma~\ref{dir:row-kernel}.
For a fixed pair of normalized exact-width sublists, write $F,G$ for
their respective total weights.

For large $\eta_1$, edges with
$\sum_{p\mid st,\ p>\eta_1^{c_1}}1/p\le1$ cost only a constant times
the product of vertex weights. On a dyad $\eta_1\asymp H_1\ge2$,
give the other edges the extra factor
\[
 \exp\left(C'\log H_1
       \sum_{\substack{p\mid st\\p\gg H_1^{c_1}}}\frac1p\right),
\]
with its cutoff no higher than required by the preceding condition.
Together with the original mismatch product, this is an admissible
inflation $h$ for \eqref{eq:lcm}, uniformly in $H_1$:
$h_p-1=O((\log p)/p)$. Since the extra factor is at least a fixed
multiple of $H_1^{C'}$ on these edges, their uninflated contribution is
\[
 O\bigl(
 H_1^{-C'+2(1-u)}
 \mathcal R^{-(1-u)}(FG)^u
 \bigr).
\]
Here $u>1/2$ is fixed.
The scale sum converges. The width sum uses geometric decay and
Cauchy--Schwarz, and the group sum uses disjointness of row assignments.
Scales with $\eta_1$ bounded away from zero are similarly covered by
\eqref{eq:lcm}.

The only remaining task is therefore to obtain the saving
\eqref{eq:small-saving} at small $\eta_1$. That saving is summable even
after multiplication by $1/\eta_1$.

\subsection{A static discard near low rational grids}

We now define the additional birth rule promised in
Proposition~\ref{dir:birth-comparisons}. Its definition uses only the
aggregate good-row lists \eqref{dir:aggregate-vertices}, not an extracted
box and not any previous table deletion.

\begin{construction}[Low rational holes]\label{dir:hole-rule}
Fix a medium group and exact width, and let $F$ be its aggregate list
mass. For each integer $\lambda\ge\lambda_0$ put $\eta=e^{-\lambda}$.
For every dyad $[A,2A)$ with $A\ge\eta^{-16}$, test all positive integer
pivots $L_*$ satisfying
\[
 \eta^4\le D_0'=r_i'L_*A^2\le\eta^{-1}
\]
whose incidence class
\[
 \{V: V/(V,L_*)\in[A,2A)\}
\]
has list mass at least $\eta^{.1}D_0'$.

For each such pivot, test only rows in its incidence class for which
$n=V/(V,L_*)$ is squarefree. On these rows, discard a birth point $x_i$ if
\[
 \left\|L_*P x_i-\frac{b}{q}\right\|\le\frac1A
 \quad\hbox{for some }1\le q\le\eta^{-.1},\quad b\in\mathbb Z.
\]
\end{construction}

Only finitely many parameters are relevant in a finite block: a
nonempty incidence class bounds $A$ by the largest low denominator, then
$A\ge e^{16\lambda}$ bounds $\lambda$, and the interval for $D_0'$
bounds $L_*$. Although evaluating the rule uses the point's phase and
masks, the families of pivots are fixed by the original aggregate lists.
In particular, this definition makes no use of a tail-prime coordinate
of the numerator.

\begin{lemma}[Cost of low rational holes]\label{dir:holes}
For every fixed $\delta>0$, Construction~\ref{dir:hole-rule} loses at
most $\delta W$ in total birth first mass when $\lambda_0$ is sufficiently
large. The choice is uniform over all groups and widths. The rule
depends on numerator residues only through squarefree low moduli.
\end{lemma}

\begin{proof}
Fix a group, width, $\lambda$, and $A$. Group the pivots by dyads, so
that $D_0'$ has a common value $D$ up to fixed factors. For two pivots
$L',L''$ in such a dyad, put
\[
 g=(L',L''),\qquad k'=L'/g,\qquad l'=L''/g.
\]
Their incidence classes have intersection mass at most
\begin{equation}\label{dir:pivot-intersection}
 C D(k'l')^{-1/2}\tau(k')\tau(l').
\end{equation}
To see this, set
\[
 u=\frac{V}{(V,g)},\qquad k=(u,k'),\qquad l=(u,l'),\qquad u=nkl.
\]
Since $(k',l')=1$, the integers $k,l$ are coprime. The two incidence
conditions imply $nk,nl<2A$. For fixed $u$, all possible $V$ divide
$gu$, so their aggregate weight is at most
$r_i'\sum_{V\mid gu}\phi(V)=r_i'gu$. For fixed $k,l$, summing $nkl$
over $n\le2A/\max(k,l)$ costs $O(A^2)$. Summing over divisors
$k\mid k'$, $l\mid l'$ gives
$O(r_i'g A^2\tau(k')\tau(l'))$, which equals the upper bound
\eqref{dir:pivot-intersection} because $L',L''$ lie in the same dyad.

If $M_0$ pivots in this dyad qualify, incidence Cauchy--Schwarz and
Lemma~\ref{ar:tensor} give, for any sufficiently small fixed
$\epsilon>0$,
\[
 (M_0\eta^{.1}D)^2
       \ll F D\,M_0^{1+\epsilon}.
\]
The divisor factors in \eqref{dir:pivot-intersection} are absorbed by
an arbitrarily small loss in its exponent price. Thus
$M_0^{1-\epsilon}D\ll F\eta^{-.2}$. Since $F\le1$ and
$D\gg\eta^4$, putting $X=M_0D$ gives
\[
 X\ll
 F^{1/(1-\epsilon)}
 \eta^{-(.2+4\epsilon)/(1-\epsilon)}
 \ll F\eta^{-1}
\]
for, say, $\epsilon<.1$. Thus
\begin{equation}\label{dir:pivot-total}
 M_0D\ll F\eta^{-1}.
\end{equation}
Each individual incidence class also has mass $O(D)$.

On an actual tested row the values of $L_*P x_i$ form a translated
uniform grid of size
$n=V/(V,L_*)\asymp A$: indeed $P x_i=m_i(z_i+\alpha_i)/V$ and
$m_i$ is a unit modulo $V$. There are $O(\eta^{-.2})$ rational
residues of denominator at most $\eta^{-.1}$. A neighborhood of
radius $1/A$ around one residue captures $O(1/A)$ of this grid.
Conditioning on the hard birth laws inflates this proportion by at
most
\[
 R(n)^{-C}\ll A^{.1}.
\]
Here $n$ is squarefree; the tag coordinates are independent of its
low digits by CRT, and double exclusions occur only at large primes,
so the conditioning estimate follows prime by prime. Dropping the
soft tests can only enlarge this loss. The hole probability relative
to hard first mass is consequently
$O(\eta^{-.2}A^{-.9})$.

Combining with \eqref{dir:pivot-total}, and weakening exponents
slightly, gives the convenient bound
$O(F\eta^{-2}A^{-.8})$ per pivot dyad. For each $A$ there are
$O(1+\lambda)$ pivot dyads, by the permitted range of $D_0'$.
The sum over $A\ge\eta^{-16}$ and then over
$\lambda\ge\lambda_0$ is an arbitrarily small multiple of $F$.
Summing the disjoint group and width masses gives the result.
The query on a row uses only the squarefree modulus $n$; hence it
has the asserted numerator dependence.
Since the hard birth mean is at least $c_{\rm hard}$ times assigned
first mass, choosing $\delta$ smaller by this fixed factor makes the
loss an arbitrarily small fraction of that hard mean as well.
\end{proof}

\subsection{Excluding saturated small-scale boxes}

We use the small-scale bins and residual edge graph defined above.
First set notation common to the three cases below. Whenever the
saving \eqref{eq:small-saving} fails and we invoke
Proposition~\ref{ar:extraction}, it gives, for an arbitrarily small
fixed $\vartheta>0$, a saturated configuration
\begin{equation}\label{dir:saturated-form}
 \begin{aligned}
 V&=NU_0a_0/b_0,&\qquad
 W'&=NT_0c_0/d_1,\\
 s&=T_0c_0b_0,&
 t&=U_0a_0d_1,
 \end{aligned}
\end{equation}
with the coprimalities of that proposition and dyads
$a_0\asymp A$, $c_0\asymp C_0$. In particular,
$\mathcal R,U_0T_0\le\eta^{-\vartheta}$,
the four varying factors are pairwise coprime and squarefree, and
\[
 D_1=r_i'NU_0A^2,\qquad D_2=r_j'NT_0C_0^2,
 \qquad
 \eta^{2+\vartheta}\le D_i\le\eta^{-\vartheta}.
\]
The captured edge weight is at least
$\eta^\vartheta D_1D_2$. On its edges,
\[
 \|s m_i\alpha_i-t m_j\alpha_j\|\ll\eta.
\]

We use the unrestricted box sampling of
Corollary~\ref{ar:box-sampling}, together with the independent mask
laws, followed on each side by the conditional anchor law of its
sampled tagged mask. An event of probability $O(\eta^{10\vartheta})$ cannot account
for this saturation. The same is true of a first-list set of mass
$O(\eta^{10\vartheta}D_1)$. All switches below are made within these
unrestricted laws; the changed values need not themselves form edges.
We retain the necessary predicate of an actual cohit, passing all
filters and not an exact primary-modeled equality. Thus for relevant
masks and anchors there are actual integral numerators
$Z_i=m_i z_i$, $Z_j=m_j z_j$ with
\begin{equation}\label{dir:actual-small-determinant}
 |s(Z_i+m_i\alpha_i)-t(Z_j+m_j\alpha_j)|\ll\eta.
\end{equation}
The outer phases depend only on their respective tail masks, not on the
box's dyadic or defect parameters. These facts specify the sampling
interface needed in each of the three cases below. In the deeper medium
case, the small-signature family \eqref{dir:small-signatures} has already
received its linear charge and is absent from this graph.
The contradictions will come from tag exclusions on high type, density
or the static holes on shallow medium type, and broad-model recognition
on deeper medium type.

\begin{lemma}[High-type boxes]\label{dir:high-box}
No saturated configuration \eqref{dir:saturated-form} occurs on high
type among the remaining non-ultra comparisons.
\end{lemma}

\begin{proof}
Here $\lambda\le Q_0Y_d+O(1)$. Since $f=1$, the tagged masks are
disjoint; by the occupancy purge each has at least $c\kappa Y_d$
prime omissions. Switch the sampled anchor on each side, using the
prime alternative in Lemma~\ref{ar:rotation} with budget
$\epsilon'\lambda$, where
$C_1\epsilon'$ is small and $\vartheta\ll\epsilon'$. Stripping the
prime from its mask costs $O(1/p)$ under the deficit law. The anchor
selection probability is at most $O(1/(\kappa Y_d))$. Therefore the
failure probability on one side is at most
\[
 \frac{C e^{-\epsilon'\lambda}}{\kappa Y_d}
 \sum_{\substack{H:\ p\asymp H\\\text{in the dense bin}}}
       \frac1{\log H}
 \ll_\kappa e^{-\epsilon'\lambda}.
\]
The sum of reciprocal logarithms over these ordinary dyads is
$O(Y_d)$. This failure is negligible under the box sampling.

Successful switches give unique coupled rationals $A_i,A_j$ for
$s m_i\alpha_i,t m_j\alpha_j$, of denominators
$e^{O(\epsilon'\lambda)}$ and errors
$O(\eta e^{O(\epsilon'\lambda)})$. By the precision and
\eqref{dir:actual-small-determinant}, their exact difference is the
negative integral determinant. Each reduced numerator is divisible
by its own sampled anchor. As in Lemma~\ref{dir:tag-fold}, the
predictions
\[
 \mathcal L_i-\frac{S_i}{[V,W']}A_i,\qquad
 \mathcal L_j-\frac{S_j}{[V,W']}A_j
\]
have sizes and denominators at most $e^{O(Y_d)}$, by
\eqref{eq:direct-size}, and fit these anchor choices. Their exact
lifted relation contradicts survival at an anchor in the disjoint
masks. Saturation therefore cannot occur.
\end{proof}

\begin{lemma}[Shallow medium boxes]\label{dir:shallow-box}
No saturated configuration occurs on medium type with
$\lambda<\ell_{\rm sh}B$ after Construction~\ref{dir:hole-rule}.
\end{lemma}

\begin{proof}
The extracted size bounds imply, for small enough $\vartheta$,
\[
 \eta^2\le A/C_0\le\eta^{-2}.
\]
\emph{Bounded dyads: the density estimate.}
First suppose both $A,C_0\le\eta^{-20}$. For a fixed $a_0$ put
$L'=PNU_0a_0$. Every first-list row at the structured vertex divides
$L'$, and the lcm and width bounds give
\[
 \eta^{23}\le L'r_i\le C\eta.
\]
In detail,
\[
 L'r_i=\eta_1\frac{r_i/r_{\max}}{T_0c_0}
       \ge c\eta^{21+2\vartheta}\ge\eta^{23}
\]
for sufficiently small $\eta$, while $T_0c_0\ge1$ gives the upper
bound.
The largest original popular host of such a row has minus-log-scale
at least its minimized cost, which is at least $B$. Since
$23\lambda<B$, the integer $L'$ is at least that original host.
The density bound \eqref{eq:centre-density} consequently applies to
the assigned structured subfamily dividing $L'$. Its density is
$O_{k_0}((L'r_i)^{H_0})$ whenever it contributes.
Its entire first mass is therefore at most
$O_{k_0}(\eta^{H_0}L'r_i)$, including all its low defects and masks.
Summing over $a_0$ uses
$\sum_{a_0\asymp A}L'r_i\asymp r_i'NU_0A^2=D_1$.
The first-list mass is consequently
\[
 O_{k_0}(\eta^{H_0}D_1),
\]
which is negligible for saturation when $\vartheta$ is sufficiently
small and $\lambda_0$ sufficiently large.

\emph{Large dyads: a qualifying static hole.}
Otherwise the ratio bound implies $A,C_0\ge\eta^{-17}$.
Switch $a_0,c_0$ as whole factors, using the integer alternative in
Lemma~\ref{ar:rotation} with budget $\epsilon'\lambda$. Its failure
probability is negligible
under the box sampling. On the remaining configurations there are
rationals $\xi_i,\xi_j$ of denominators
$e^{O(\epsilon'\lambda)}$ satisfying
\[
 \left|T_0b_0m_i\alpha_i-\xi_i\right|
       \ll\frac{\eta e^{O(\epsilon'\lambda)}}{C_0},
 \qquad
 \left|U_0d_1m_j\alpha_j-\xi_j\right|
       \ll\frac{\eta e^{O(\epsilon'\lambda)}}{A}.
\]
Rational separation and \eqref{dir:actual-small-determinant} force
\begin{equation}\label{dir:common-rational}
 c_0(T_0b_0Z_i+\xi_i)
       =a_0(U_0d_1Z_j+\xi_j).
\end{equation}
Set $L_*=NU_0T_0$ and $y_i=P x_i$. The common rational
\[
 \frac{T_0b_0Z_i+\xi_i}{a_0}
       =\frac{U_0d_1Z_j+\xi_j}{c_0}
\]
has denominator at most $e^{O(\epsilon'\lambda)}$: its denominator
divides both $a_0$ times the denominator of $\xi_i$ and $c_0$ times
that of $\xi_j$, and $(a_0,c_0)=1$. Moreover,
$L_*y_i$ is within
$O(\eta e^{O(\epsilon'\lambda)}/(AC_0))$ of that rational.
The quotient
\[
 \frac{V}{(V,L_*)}=a_0
\]
is squarefree, even when $a_0$ shares primes with $N$; its additional
prime exponents are precisely the ones surviving this quotient.

This pivot qualifies in Construction~\ref{dir:hole-rule}.
Indeed $D_0'=T_0D_1$ lies between $\eta^4$ and $\eta^{-1}$ for small
$\vartheta$. If its incidence class had mass less than
$\eta^{.1}D_0'$, it would have mass at most
$\eta^{.1-\vartheta}D_1\le\eta^{10\vartheta}D_1$ when
$\vartheta<.1/11$. This is a restriction on aggregate vertex mass,
already integrated over the auxiliary masks and anchors, so
Corollary~\ref{ar:box-sampling} contradicts saturation. Also
$A\ge\eta^{-17}$ lies in the tested
range. Taking $\epsilon'$ sufficiently small makes the rational
denominator at most $\eta^{-.1}$ and the approximation error less
than $1/A$. Thus the actual point fails its hole test, contradicting
the retained cohit predicate.
\end{proof}

At this stage high boxes and shallow medium boxes have been excluded.
The remaining medium range has $\lambda$ comparable to $B$. This
allows small relative phase errors to activate a broad tier. The
preliminary charge has removed the case where all signature factors
and the width ratio are too small for a switch.

\begin{lemma}[Deeper medium boxes]\label{dir:deep-box}
For medium type with
$\ell_{\rm sh}B\le\lambda\le Q_0B+O(1)$, an additional nonexact
comparison count of size $O(W)$ may be charged directly. No
saturated configuration remains after that charge.
\end{lemma}

\begin{proof}
The $O(W)$ charge for \eqref{dir:small-signatures} was established
before defining the aggregate graph. It remains to exclude saturation
among its other edges.

\emph{The switch families.}
In a resulting saturated box, switch the following factors on side
$i$: the whole factor $c_0$ when $C_0\ge\eta^{-t_1}$; every divisor
of $b_0$ in dyads between $2\eta^{-t_1}$ and $\eta^{-J_1}$; and
every prime of $m_i$. On side $j$ use $a_0,d_1,m_j$.
Here $t_1>0$ is sufficiently small against $h'/Q_0$,
$J_1\gg J/\ell_{\rm sh}$ is fixed, then
$C_1\epsilon'$ is small against $t_1$, and finally $\vartheta$ is
sufficiently small.

For a divisor $h$ of the squarefree $b_0$, stripping under $\pi_N$
costs at most
\[
 \frac1{\phi(h)}
       \ll\frac{e^{\epsilon'\lambda/4}}h
       \qquad(h\le\eta^{-J_1}).
\]
At a prime $p\mid h$, the ratio of the single-omission law to the
zero-omission law is $1/p$ if $\nu_p(N)\ge2$, and $1/(p-1)$ if
$\nu_p(N)=1$. Their product is at most $1/\phi(h)$.
The ratio is applied only to valid original squarefree configurations;
dropping restrictions on the stripped parameters then enlarges the
upper count. No switched configuration is required to be good or
assigned.
The integer rotation estimate therefore makes the union of divisor
switch failures negligible, including the $O_{J_1}(\lambda)$ dyads.
Whole-factor switches have the same saving. A prime stripped from
$m_i$ costs $O(1/p)$; the prime rotation estimate, summed through
log-log heights $O(B)$, loses at most an additional factor $O(1+B)$.
Since $\lambda\ge\ell_{\rm sh}B$, these prime failures are also
negligible under the box sampling.

\emph{Qualification of the broad model.}
Outside all failures, any switch gives unique coupled approximants
to $s m_i\alpha_i,t m_j\alpha_j$, of denominators
$e^{O(\epsilon'\lambda)}$, errors
$O(\eta e^{O(\epsilon'\lambda)})$, and exact difference equal to
the negative determinant. After division by the whole signatures,
their denominators are at most $e^{O(t_1)\lambda}$. Here are all
the factors in this assertion. Every prime of $m_i$ divides the
first approximant numerator. Division by a switched $c_0$ leaves
a small denominator; if it is not switched, $c_0<2\eta^{-t_1}$.
Among prime factors of $b_0$ not dividing the numerator, the product
cannot reach the switch range: all are at most $e^{JB}$, so a
first-crossing product is still below $\eta^{-J_1}$ and would
contradict the small denominator after its switch. These factors
are coprime, and the additional factor $T_0$ is at most
$\eta^{-\vartheta}$. The second side is identical.

The errors after division, measured in $\psi_i,\psi_j$, are at
most $e^{O(\epsilon'+\vartheta)\lambda}$, using the extracted
width-ratio bound. Because $\lambda\le Q_0B+O(1)$ and the small
parameters were chosen in the stated order, the denominators and
relative heights fit the applicable broad tier. That tier is
enabled by $\psi_i,\psi_j\ll e^{-\ell_{\rm sh}B}$.
The exact coupled relation now makes the cohit an exact
primary-modeled equality, again a contradiction.

If no switch is available on either side, the same factor bounds
give $s m_i,t m_j\le e^{O(t_1)\lambda}$; the extracted width ratio
is at most $\eta^{-\vartheta}$. These fit the small-signature and
width bounds of the already charged family. This too is excluded,
and hence saturation is impossible.
\end{proof}

\begin{proof}[Proof of Proposition~\ref{dir:birth-comparisons}]
Lemma~\ref{dir:holes} supplies the static birth discard, its loss,
and its squarefree low-coordinate dependence.
Lemma~\ref{dir:ultra-count} gives the asserted narrow estimate and
the same estimate for the ultra broad portion. For the remaining
broad portion, exact primary-modeled equalities have a linear
charge by Lemma~\ref{dir:exact-equalities}. The preliminary charge
for \eqref{dir:small-signatures} is also $O(W)$, and
Lemma~\ref{dir:row-kernel} reduces the residual counts to weighted
arithmetic edges.

The large and bounded-away-from-zero determinant scales were summed
above. At small scales, Lemmas~\ref{dir:high-box},
\ref{dir:shallow-box}, and \ref{dir:deep-box} exclude every
saturation alternative in this residual graph. Thus
\eqref{eq:small-saving} holds on all remaining edges. After the
factor $1/\eta_1$ in \eqref{eq:row-kernel}, its scale contribution
is
\[
 O\bigl(\eta^\chi\mathcal R^{-1/10}(FG)^u\bigr).
\]
The $\lambda$ sum converges; summing dyadic widths uses the
geometric factor $\mathcal R^{-1/10}$ and Cauchy--Schwarz, and
summing groups uses disjointness. Since $W\le1$ and $2u>1$, these
terms are bounded independently of the block. All steps used
unclipped birth weights, so none of the constants depends on
$N_0$. Finally, further thinning decreases each nonnegative
comparison count. This proves all assertions.
\end{proof}
